\caption{Operational uncertainty, rather than semantic proximity alone, determines when archived causal evidence can be reused. \textbf{A}: although every recorded descriptor is identical, a hidden mechanism shift makes the sampling-only interval undercover; the robust interval retains at least nominal coverage throughout its declared uncertainty set (shading). Ribbons are 95\% Wilson Monte Carlo intervals. \textbf{B}: selecting the largest absolute coordinate makes pointwise coverage collapse with archive size, while simultaneous Bonferroni protection remains valid. Curves are exact Gaussian probabilities and markers are Monte Carlo estimates. \textbf{C}: in the anchored Gaussian model, the empirical worst-case mean absolute error follows $\max\{Ld,I^{-1/2}\}$, with the dashed line separating ambiguity- and noise-dominated regimes; white curves are iso-risk contours. \textbf{D}: for linear-Gaussian bridge design, log-determinant greedy selection nearly matches exhaustive search across 200 generated design instances and exceeds the worst-case $1-1/e$ guarantee. The inset gives an exact non-submodularity witness for partial-identification (PI) width: for $\beta\in[-1,1]^2$ and target $\beta_1$, bridge $A$ observes $\beta_1+\beta_2=0$ and bridge $B$ observes $\beta_2=0$; neither reduces width alone, whereas their combination point-identifies the target. All panels are synthetic, with seeds and replication counts in Table~\ref{tab:operational-simulation-manifest}.}
